\section{A measured residual fibre for the marked quaternionic cubic}
\label{sec:residual-measure}

The measured-fibre integration format in Nolan's \emph{Measured-Bundle
Projection Framework}, Sections 2.2--3.1,
\cite{Nolan2025Threads,NolanRepository2025MBPF} suggests applying conditional
integration to the compact residual fibre of the previously constructed
nonlinear metric map. This section carries out that application. It retains
the original fifteen-dimensional matrix, its determinant-one marking, its
Moore cubic, and all four metric weights. The resulting conditional measure,
cubic probability law, positive-cone restriction, and weighted moments are
computed explicitly. The construction uses the earlier metric map with its
documented provenance; the integration and formulas below are an additional
application of the cited measured-bundle idea.

\subsection{Original coordinates, cubic, and determinant-one marking}

Fix $\mathbb C=\mathbb R+\mathbb R\mathbf i\subset\mathbb H$,
$\mathbf j^2=-1$, and $\mathbf ja=\bar a\mathbf j$ for $a\in\mathbb C$.
Quaternionic conjugate transpose is $*$. Keep the matrices
\[
 P=\begin{pmatrix}1&-1&0\\0&1&-1\end{pmatrix},\quad
 u=\begin{pmatrix}1\\1\\1\end{pmatrix},\quad
 C=PP^*=\begin{pmatrix}2&-1\\-1&2\end{pmatrix},
\]
\[
 R=P^*C^{-1}=\frac13\begin{pmatrix}2&1\\-1&1\\-1&-2\end{pmatrix},\quad
 X=(R\;u),\quad
 X^{-1}=\begin{pmatrix}P\\u^*/3\end{pmatrix},\quad
 X^*X=\operatorname{diag}(C^{-1},3).
\]
Multiplication verifies these identities and $\det_{\mathbb R}X=1$.
For $H\in J=\operatorname{Herm}_3(\mathbb H)$ write
\begin{equation}\label{rm:marked-matrix}
 H=XMX^*,\qquad
 M=\begin{pmatrix}K&q\\q^*&\lambda\end{pmatrix},\quad
 K=G+\bar\rho\varepsilon\mathbf j,\quad q=z+w\mathbf j,
 \quad\varepsilon=\begin{pmatrix}0&1\\-1&0\end{pmatrix},
\end{equation}
where
\[
 G=\begin{pmatrix}a&c\\\bar c&b\end{pmatrix},\quad
 a,b,\lambda\in\mathbb R,\quad c,\rho\in\mathbb C,\quad
 z,w\in\mathbb C^2,\quad S=[z,w].
\]
Every displayed coordinate is recovered uniquely from $X^{-1}HX^{-*}$;
the real coordinate count is $4+4+4+2+1=15$.

For quaternionic Hermitian matrices the cubic used throughout is
\begin{equation}\label{rm:moore-polynomial}
 N\!\begin{pmatrix}a&v&q_1\\\bar v&b&q_2\\
 \bar q_1&\bar q_2&\lambda\end{pmatrix}
 =ab\lambda-\lambda|v|^2-b|q_1|^2-a|q_2|^2
       +2\operatorname{Re}(\bar q_1vq_2).
\end{equation}
This convention retains the cyclic order in the last term.
The marking satisfies $N(XMX^*)=N(M)$. A verification that does not
change the marking is as follows. For a quaternionic matrix $A+B\mathbf j$,
the complex representation
\[
 \mathcal R(A+B\mathbf j)
 =\begin{pmatrix}A&B\\-\bar B&\bar A\end{pmatrix}
\]
preserves products and adjoints, by the multiplication rule
\begin{equation}\label{rm:quaternion-product}
 (a_0+b_0\mathbf j)(c_0+d_0\mathbf j)
 =(a_0c_0-b_0\bar d_0)+(a_0d_0+b_0\bar c_0)\mathbf j.
\end{equation}
For a Hermitian $2\times2$ quaternionic matrix
$\left(\begin{smallmatrix}b&v\\\bar v&d\end{smallmatrix}\right)$,
complex Schur elimination, first for $b\ne0$, gives complex determinant
$(bd-|v|^2)^2$. Polynomial continuation proves this also for $b=0$.
Eliminating the first row and column of a Hermitian $3\times3$ matrix
with $a\ne0$ gives
\[
 a\left[(b-|v|^2/a)(\lambda-|q_1|^2/a)
             -|q_2-\bar vq_1/a|^2\right]=N(M).
\]
Expansion verifies the equality with \eqref{rm:moore-polynomial}; the
terms $|v|^2|q_1|^2/a$ cancel. Applying the preceding complex determinant
calculation to this Schur complement yields
$\det_{\mathbb C}\mathcal R(M)=N(M)^2$, again extended by polynomial
continuation. Now $\det_{\mathbb C}\mathcal R(X)=(\det X)^2=1$, so
$N(XMX^*)^2=N(M)^2$. In the integral domain of real polynomials in the
fifteen entries, this forces one of the polynomial factors
$N(XMX^*)-N(M)$ and $N(XMX^*)+N(M)$ to vanish identically. At $M=I_3$,
the ordinary real symmetric determinant formula gives
$N(XX^*)=\det(XX^*)=1=N(I_3)$; the plus factor does not vanish there.
The minus factor is therefore the zero polynomial, as asserted.

\subsection{The metric map and its entire compact fibre}

Use the fixed Pauli basis
\[
 \sigma_0=I_2,\quad \sigma_1=\begin{pmatrix}0&1\\1&0\end{pmatrix},
 \quad\sigma_2=\begin{pmatrix}0&-i\\i&0\end{pmatrix},
 \quad\sigma_3=\begin{pmatrix}1&0\\0&-1\end{pmatrix},
 \quad\eta=\operatorname{diag}(1,-1,-1,-1).
\]
Thus $Y=y^a\sigma_a$ has determinant $y^T\eta y$, with positive
definite $Y$ corresponding to the future timelike cone. For invertible
$S$ define, without discarding its scale,
\begin{equation}\label{rm:metric-map}
 \begin{gathered}
 r=|\det S|,\qquad L_S(Y)=SYS^*/r,\\
 (d_0,d_1,d_2,d_3)=(r,e^{\operatorname{Re}\rho},
                         e^{\operatorname{Im}\rho},e^\lambda),\\
 h=L_S^{-T}\eta\operatorname{diag}(d_0,d_1,d_2,d_3)L_S^{-1}.
 \end{gathered}
\end{equation}
The source domain $\mathcal U$ consists of the matrices
\eqref{rm:marked-matrix} for which $\det S\ne0$, the four $d_a$ are
pairwise distinct, and $d_1<d_2<d_3$. Put $\Psi(H)=(G,h)$.
The target is $\mathcal B=\operatorname{Herm}_2(\mathbb C)\times\mathcal C$,
where $\mathcal C$ consists of forms whose relative endomorphism
$B_h=\eta^{-1}h$ has positive distinct eigenvalues and a future oriented,
positively oriented $\eta$-orthonormal eigenframe, with increasing spatial
eigenvalue labels. The metric with the opposite signature convention is
$-h$ with background $-\eta$; its relative endomorphism is exactly the
same $B_h$. All formulas below therefore apply unchanged to $(G,-h)$
with that background.

Define
\[
 Q_8=\{\pm I_2,\ \pm i\sigma_1,\ \pm i\sigma_2,\ \pm i\sigma_3\},
 \quad\mathcal R_0=\{I_2,i\sigma_1,i\sigma_2,i\sigma_3\},
\]
\begin{equation}\label{rm:compact-group}
 F=\{e^{i\theta}q_0:q_0\in Q_8,\ \theta\in\mathbb R\}
   \simeq (U(1)\times Q_8)/\{(1,I_2),(-1,-I_2)\}.
\end{equation}
It acts on the original matrices through
$(G,S,\rho,\lambda)\mapsto(G,SU,\rho,\lambda)$, $U\in F$,
and then \eqref{rm:marked-matrix}. This is a specified action on $J$,
with $X$ fixed.

\begin{proposition}\label{rm:fibre}
The map $\Psi:\mathcal U\to\mathcal B$ is a smooth surjective submersion
of rank fourteen. Its fibres are exactly the free $F$-orbits. Each orbit
has four circle components, and local smooth sections of $\Psi$ exist.
For $y=(G,h)$, let $v_0(y)$ be the unique future unit timelike eigenvector
of $B_h$, and put $V_0(y)=\sum_{a=0}^3v_0^a(y)\sigma_a$. The functions
\begin{equation}\label{rm:target-recovery}
 r=d_0,\quad\rho=\log d_1+i\log d_2,\quad
 \lambda=\log d_3,\qquad T:=SS^*=rV_0(y)
\end{equation}
are independent of the point chosen in the fibre and are smooth on
$\mathcal B$.
\end{proposition}
\begin{proof}
The determinant and positivity identities
$\det(SYS^*/r)=\det Y$ and $SYS^*/r>0$ for $Y>0$ show that $L_S$
preserves $\eta$ and its future cone. The group $GL_2(\mathbb C)$ is
connected: polar decomposition connects its positive factor to $I_2$
by a straight positive path, and diagonalizing the unitary factor
connects its two eigenvalue phases to zero. Continuity then gives
$\det_{\mathbb R}L_S=1$. Direct multiplication gives $L_{ST}=L_SL_T$.
If $L_S$ is the identity, evaluating at $I_2$ gives $SS^*=rI_2$;
the unitary $S/\sqrt r$ commutes with every Hermitian matrix. Commuting
with $\sigma_3$ and $\sigma_1$ makes it scalar, so the kernel consists
of the nonzero complex scalars. Its intersection with fixed determinant
magnitude $r$ is exactly $\{\sqrt r\,e^{i\theta}I_2:\theta\in\mathbb R\}$;
the scalar phase subgroup acts transitively on this circle.

For completeness every proper future Lorentz frame has a complex spin
lift. If its future first column is $v$, the matrix
$V=\sum v^a\sigma_a$ is positive with determinant one. Its positive
square root $B$ has determinant one and $L_B(I_2)=V$. The remaining
frame after applying $L_B^{-1}$ fixes the time axis, hence is an
oriented rotation on the three spatial coordinates. An oriented
orthogonal $3\times3$ matrix has eigenvalue $1$ (its other nonreal
eigenvalues pair conjugately, and its determinant is $1$); on the
orthogonal plane it is a two-dimensional rotation. For its unit axis
$n$ and oriented angle $t$, direct Pauli multiplication gives the lift
$\cos(t/2)I_2-i\sin(t/2)n\cdot\sigma$. Multiplying this lift by $B$
gives the required $SL_2(\mathbb C)$ lift. The kernel there is
$\{I_2,-I_2\}$ by the scalar calculation and determinant one.

The identity $B_h=L_S\operatorname{diag}(d_a)L_S^{-1}$ proves the
claimed image inclusion. Conversely choose a frame $L$ and labelled
eigenvalues for any member of $\mathcal C$, lift $L$ to $g\in
SL_2(\mathbb C)$ as just proved, and set
$S=\sqrt{d_0}g$, $\rho=\log d_1+i\log d_2$,
$\lambda=\log d_3$. These coordinates give precisely the required
metric and retain $|\det S|=d_0$.

At $S=I_2$, a tangent is uniquely $A+(t+i\theta)I_2$ with
$\operatorname{tr}A=0$, $t,\theta\in\mathbb R$. The derivative of
$r$ is $2t$, and that of $L_S$ is $Y\mapsto AY+YA^*$.
If the latter is zero, evaluating at $I_2$ and then at all Hermitian
$Y$ makes $A$ scalar skew-Hermitian; its trace makes it zero.
The six real traceless complex directions therefore map isomorphically
to the six Lorentz tangent directions, with the phase as the only
remaining kernel. Translation gives this result at every $S$.
Writing $D=\operatorname{diag}(d_a)$ and a frame variation as
$L\exp(tA)$, the transported metric derivative is
\[
 L^T(\delta h)L=\eta\delta D+\eta(AD-DA).
\]
The four diagonal coordinates vary independently. For $a<b$ the
off-diagonal coordinate is $\eta_a(d_b-d_a)A_{ab}$, and the six
factors $d_b-d_a$ are nonzero. This gives ten metric directions plus
the four independent coordinates of $G$. The inverse function theorem
on a transverse fourteen-dimensional slice gives local sections and
openness of $\mathcal C$.

For two points with the same metric the labelled distinct eigenvalues
fix every $d_a$. The two Lorentz frames differ by a diagonal sign
matrix $\operatorname{diag}(1,s_1,s_2,s_3)$ with $s_1s_2s_3=1$:
indeed their quotient commutes with the diagonal simple-spectrum $D$,
and time orientation fixes the first sign. These four matrices are
exactly the Lorentz images of $Q_8$, as direct conjugation of the Pauli
matrices shows. The kernel calculation then gives $S'=Se^{i\theta}q_0$.
Conversely every such multiplication fixes the metric. Since $S$ is
invertible the action is free. The group in \eqref{rm:compact-group}
has four circle components, represented by $\mathcal R_0$.
Finally $L_S(I_2)=SS^*/r$ is the unique future unit timelike eigenvector
matrix. Simple eigenvectors vary smoothly locally by the same
inverse-function calculation, and uniqueness of the future choice
glues these vectors globally. The labelled eigenvalues and their
logarithms are smooth for the same reason. This proves
\eqref{rm:target-recovery}.
\end{proof}

\subsection{A canonical conditional probability and its exact projection}

For a Borel function $f$ on $F$ the normalized Haar integral is
\begin{equation}\label{rm:haar}
 \int_F f(U)\,d\alpha(U)
 =\frac1{8\pi}\sum_{q_0\in\mathcal R_0}
                 \int_0^{2\pi}f(e^{i\theta}q_0)\,d\theta.
\end{equation}
This formula follows by pushing forward the uniform measure on
$U(1)\times Q_8$ and pairing $q_0$ with $-q_0$ using the change
$\theta\mapsto\theta+\pi$. Its total mass is one. Multiplication in
$F$ permutes the four cosets and translates or changes the chosen
representative phase, so the displayed integral is both left and right
invariant. In particular no branch receives an unrecorded multiplicity.

For $y\in\mathcal B$ and any $H_y\in\Psi^{-1}(y)$ define a probability
measure on $\mathcal U$ by
\begin{equation}\label{rm:conditional-kernel}
 \kappa_y(E)=\int_F\mathbf1_E(H_y\cdot U)\,d\alpha(U),\qquad
 \mathcal A f(y)=\int_{\mathcal U}f(H)\,d\kappa_y(H).
\end{equation}
Here $H_y\cdot U$ denotes the action in \eqref{rm:marked-matrix}, not
ordinary matrix multiplication of $H_y$ by a $2\times2$ matrix.
If $H_y$ is changed to $H_y\cdot U_0$, left invariance of
\eqref{rm:haar} leaves \eqref{rm:conditional-kernel} unchanged. Thus
no global frame or spin lift has been chosen. The kernel is supported
on the exact fibre. A local section with column $S_y$ gives product
coordinates $(y,U)\mapsto H_y\cdot U$, whose inverse group coordinate
is the explicit smooth matrix $U=S_y^{-1}S$. The fibre proposition
places this inverse in $F$. On these local product coordinates, integration in the
fixed formula \eqref{rm:haar} makes $y\mapsto\kappa_y(E)$ Borel:
this follows first for indicators of rectangles in local product
coordinates, then for all Borel sets by closure under complements and
countable disjoint unions. A countable cover by such charts exists
because $\mathcal B$ is an open subset of a finite-dimensional
Euclidean space. Disjointizing this cover proves the assertion
globally. The same local formula and compactness of $F$ show that
$\mathcal A f$ is continuous for continuous $f$; on a small compact
target neighbourhood the integrand is uniformly continuous on the
product with $F$.

The kernel also retains the exact Lorentz equivariance of the source
map. For $g\in SL_2(\mathbb C)$ define
\[
 \mathcal T_g(H)=X\operatorname{diag}(g,1)X^{-1}
           H X^{-*}\operatorname{diag}(g,1)^*X^*,\qquad
 \ell_g(G,h)=(gGg^*,L_g^{-T}hL_g^{-1}).
\]
Since $g\varepsilon g^T=\varepsilon$ and
$\mathbf jg^*=g^T\mathbf j$, its marked coordinates are exactly
$(gGg^*,gS,\rho,\lambda)$. Consequently
$\Psi\mathcal T_g=\ell_g\Psi$ by $L_{gS}=L_gL_S$.
Left multiplication of $S$ by $g$ commutes with its right
multiplication by $F$, so \eqref{rm:conditional-kernel} gives
\begin{equation}\label{rm:kernel-equivariance}
 (\mathcal T_g)_*\kappa_y=\kappa_{\ell_g y},\qquad
 \mathcal A(f\circ\mathcal T_g)(y)=\mathcal A f(\ell_g y).
\end{equation}
This assertion concerns the explicit conditional kernels. The finite
coordinate-weighted base measure introduced next is separately
specified and is not asserted to be Lorentz invariant.

The following cubic calculation will also provide an explicit finite
base measure for a Hilbert-space realization of this kernel. Define
\begin{equation}\label{rm:mean-radius-definitions}
 \Delta=\det G-|\rho|^2,\qquad
 m(y)=\lambda\Delta-\operatorname{tr}((\operatorname{adj}G)T),
 \qquad\beta(y)=2|\rho|r,
\end{equation}
using the recovered coordinates \eqref{rm:target-recovery}. These are
smooth real functions and $\beta>0$ everywhere on $\mathcal B$:
$r>0$, while $\rho=0$ would imply $d_1=d_2=1$, excluded by the domain.
Here $\operatorname{adj}G=\left(\begin{smallmatrix}b&-c\\-\bar c&a
\end{smallmatrix}\right)$, with no division by $\det G$.

Choose the actual coordinate vector
\[
 \xi(y)=(a,b,\operatorname{Re}c,\operatorname{Im}c,
                         (h_{ab})_{0\leq a\leq b\leq3})\in\mathbb R^{14}
\]
in a fixed lexicographic order of the last ten entries, and define
\begin{equation}\label{rm:concrete-measures}
 \begin{gathered}
 Z=\int_{\mathcal B}\frac{e^{-|\xi(y)|^2}}{1+m(y)^2+\beta(y)^2}\,d^{14}\xi,
 \\d\mu(y)=Z^{-1}
 \frac{e^{-|\xi(y)|^2}}{1+m(y)^2+\beta(y)^2}\,d^{14}\xi,\qquad
 \nu(E)=\int_{\mathcal B}\kappa_y(E)\,d\mu(y).
 \end{gathered}
\end{equation}
These are specified additional probability measures; the original
matrix and metric have not been rescaled. The constant $Z$ is retained.
It is positive because $\mathcal B$ is nonempty and open (for example
$S=\operatorname{diag}(2,1)$ and weights $(2,3,4,5)$ supply a point).
It is finite because its integrand is bounded by the integrable
fourteen-dimensional Gaussian. Countable additivity of $\nu$ follows
from countable additivity of each $\kappa_y$ and monotone convergence;
$\nu(\mathcal U)=1$. The support property gives
\begin{equation}\label{rm:pushforward}
 \Psi_*\nu=\mu,\qquad
 \int_{\mathcal U} f(H)\,d\nu(H)
 =\int_{\mathcal B}\left(\int_{\mathcal U}f(H)\,d\kappa_y(H)\right)d\mu(y)
\end{equation}
first for indicators, then simple functions, then nonnegative
functions by monotone convergence, and then absolutely integrable
functions by taking positive and negative parts.

\begin{proposition}\label{rm:hilbert-projection}
The operator $\mathcal A:L^2(\mathcal U,\nu)\to L^2(\mathcal B,\mu)$
has norm one. Pullback $\Psi^*$ is an isometry,
$\mathcal A\Psi^*=I$, and $\mathcal P=\Psi^*\mathcal A$ is the
orthogonal projection onto the closed subspace of fibre-constant
$L^2$ functions. Explicitly,
\begin{equation}\label{rm:l2-remainder}
 \|f\|_{L^2(\nu)}^2
 =\|\mathcal A f\|_{L^2(\mu)}^2
 +\int_{\mathcal B}\int_{\mathcal U}
       |f(H)-\mathcal A f(y)|^2\,d\kappa_y(H)\,d\mu(y).
\end{equation}
Thus \eqref{rm:conditional-kernel} is a concrete disintegration and
conditional-expectation map for the measures \eqref{rm:concrete-measures}.
\end{proposition}
\begin{proof}
If $f=0$ $\nu$-almost everywhere, \eqref{rm:pushforward} applied to
$|f|^2$ gives $\int|f|^2\,d\kappa_y=0$ for $\mu$-almost every $y$.
The fibre Cauchy--Schwarz inequality then makes $\mathcal A f=0$
$\mu$-almost everywhere. Thus $\mathcal A$ is well defined on the
$L^2$ equivalence classes in the statement.
Cauchy--Schwarz against the probability $\kappa_y$ gives
$|\mathcal A f(y)|^2\leq\int|f|^2\,d\kappa_y$. Integrating and
using \eqref{rm:pushforward} proves the contraction; $f=1$ proves norm
one. The pushforward identity also proves that $\Psi^*$ is an
isometry. A function $g\circ\Psi$ is constant with value $g(y)$ on
the support of $\kappa_y$, so $\mathcal A\Psi^*g=g$.
For bounded $f,g$, \eqref{rm:pushforward} gives
\[
 \langle f,\Psi^*g\rangle_{L^2(\nu)}
 =\langle\mathcal A f,g\rangle_{L^2(\mu)}.
\]
Truncation and Cauchy--Schwarz extend this identity to all $L^2$
functions. Thus $\mathcal A$ is the adjoint of $\Psi^*$ and
$\mathcal P$ is self-adjoint and idempotent. The image of an isometry
from the complete space $L^2(\mu)$ is closed, so this is the asserted
orthogonal projection. Here the fibre-constant subspace means the
pullback range, with functions identified almost everywhere. A Borel
function literally constant on fibres descends to a Borel function
by evaluating on each of the local smooth sections, so this also
agrees with the pointwise meaning wherever a representative is fixed.
Alternatively expand the square inside
\eqref{rm:l2-remainder}; the cross term is
$2|\mathcal A f(y)|^2$ and the constant term is
$|\mathcal A f(y)|^2$, proving the formula directly.
\end{proof}

\subsection{The exact cubic law on all four fibre components}

\begin{theorem}\label{rm:cubic-law}
For every $H\in\mathcal U$, put $\delta=\det S$. The original cubic is
\begin{equation}\label{rm:cubic-coordinate-formula}
 N(H)=\lambda(\det G-|\rho|^2)
       -\operatorname{tr}((\operatorname{adj}G)SS^*)
       -2\operatorname{Re}(\rho\det S).
\end{equation}
On its entire fibre, with $q_0\in\mathcal R_0$ and $0\leq\theta<2\pi$,
\begin{equation}\label{rm:cubic-fourier}
 N(H\cdot(e^{i\theta}q_0))
   =m(y)-\rho\delta e^{2i\theta}
          -\overline{\rho\delta}\,e^{-2i\theta},\qquad y=\Psi(H).
\end{equation}
All four discrete components have the same cubic distribution. Their
common, hence full conditional, law is the probability measure
\begin{equation}\label{rm:arcsine-law}
 N_*\kappa_y(dt)
 =\frac{\mathbf1_{(m(y)-\beta(y),m(y)+\beta(y))}(t)}
        {\pi\sqrt{\beta(y)^2-(t-m(y))^2}}\,dt.
\end{equation}
In particular
\begin{equation}\label{rm:first-two-moments}
 \mathcal A N=m,\qquad
 \mathcal A(N^2)=m^2+2|\rho|^2r^2,\qquad
 \operatorname{Var}_{\kappa_y}(N)=2|\rho|^2r^2=\beta^2/2>0.
\end{equation}
For every integer $k\geq0$ the complete moment formula is
\begin{equation}\label{rm:all-moments}
 \mathcal A(N^k)(y)=
 \sum_{j=0}^{\lfloor k/2\rfloor}
  \binom{k}{2j}\binom{2j}{j}
       m(y)^{k-2j}|\rho\delta|^{2j}.
\end{equation}
The probability law, mean, and variance are target observables although
$N$ itself is not constant on any regular fibre.
\end{theorem}
\begin{proof}
Write $q_j=z_j+w_j\mathbf j$ and $v=c+\bar\rho\mathbf j$ in
\eqref{rm:moore-polynomial}. Applying \eqref{rm:quaternion-product}
twice gives
\[
 \operatorname{Re}(\bar q_1vq_2)
 =\operatorname{Re}\bigl(
  \bar z_1cz_2+w_1\rho z_2
        -\bar z_1\bar\rho\bar w_2+w_1\bar c\bar w_2\bigr)
 =\operatorname{Re}\bigl(c(\bar z_1z_2+\bar w_1w_2)-\rho\delta\bigr).
\]
Here taking complex conjugates inside real parts changes the third
term to $-\operatorname{Re}(\rho z_1w_2)$ and the fourth to
$\operatorname{Re}(c\bar w_1w_2)$; this fixes both the sign and the
order of the determinant term. Also $|v|^2=|c|^2+|\rho|^2$ and
$|q_j|^2=|z_j|^2+|w_j|^2$. Expanding the trace of
$(\operatorname{adj}G)SS^*$ gives
\[
 b(|z_1|^2+|w_1|^2)+a(|z_2|^2+|w_2|^2)
             -2\operatorname{Re}(c(\bar z_1z_2+\bar w_1w_2)),
\]
which proves \eqref{rm:cubic-coordinate-formula}.

Every $q_0\in Q_8$ has complex determinant one and is unitary.
Consequently $S'=Se^{i\theta}q_0$ has
$S'S'^*=SS^*$ and $\det S'=e^{2i\theta}\delta$. Substitution proves
\eqref{rm:cubic-fourier}. The discrete invariance also has an exact
quaternionic congruence interpretation. Every matrix
$k=\left(\begin{smallmatrix}\alpha&\gamma\\-\bar\gamma&\bar\alpha
\end{smallmatrix}\right)\in SU(2)$ corresponds to the unit quaternion
$t=\alpha+\gamma\mathbf j$. Right multiplication $q\mapsto qt$
changes $[z,w]$ to $[z,w]k$ by \eqref{rm:quaternion-product}.
It is the congruence by $\operatorname{diag}(I_2,\bar t)$ on $M$,
transported through the unchanged $X$. In particular the discrete
components do not change positivity either.

Set $\rho\delta=|\rho\delta|e^{i\varphi}$. Under the uniform phase,
$2\theta+\varphi$ modulo $2\pi$ is uniform: splitting its integral
into two intervals of length $\pi$ proves this by substitution.
Thus $N$ has the law of $m-\beta\cos t$ for uniform $t$.
On $(0,\pi)$ this function is increasing with derivative
$\beta\sin t=\sqrt{\beta^2-(N-m)^2}$; on $(\pi,2\pi)$ it is
decreasing with the opposite derivative. The two changes of variables
give \eqref{rm:arcsine-law}, including its normalization. The endpoints
have probability zero. Integration of the nonzero Fourier modes in
\eqref{rm:cubic-fourier} gives mean $m$; expanding its square gives
the constant term $m^2+2|\rho\delta|^2$. For general $k$, the only
terms surviving integration contain the same number $j$ of the
$-\rho\delta e^{2i\theta}$ and the conjugate negative-frequency
factor. Their number is $\binom{k}{2j}\binom{2j}{j}$ and their sign
is positive, proving \eqref{rm:all-moments}. Finally $\beta>0$ was
proved after \eqref{rm:mean-radius-definitions}, so every such law
has strictly positive variance.
\end{proof}

The concrete measure in \eqref{rm:concrete-measures} makes this cubic
an actual $L^2(\nu)$ function: \eqref{rm:first-two-moments} and
\eqref{rm:pushforward} give
\[
 \int N^2\,d\nu
 =Z^{-1}\int_{\mathcal B}
 \frac{m^2+\beta^2/2}{1+m^2+\beta^2}e^{-|\xi|^2}\,d^{14}\xi<\infty.
\]
Thus the Hilbert-space projection statement applies to this specific
original observable, rather than only to bounded test functions.
Its nonzero discarded squared norm is exactly
$\int_{\mathcal B}2|\rho|^2r^2\,d\mu$.

For a fixed real number $c_0$ and retained coefficient $\zeta\geq0$,
an immediate identity for a cubic penalty is
\begin{equation}\label{rm:cubic-penalty}
 \mathcal A\!\left[\frac{\zeta}{2}(N-c_0)^2\right]
 =\frac{\zeta}{2}(m-c_0)^2+\zeta|\rho|^2r^2.
\end{equation}
The last term is exactly the fluctuation contribution. Replacing
$N$ by its mean before taking its square omits it. This is an identity
of observables; no equation of motion has been changed or asserted
to intertwine by this integration.

\subsection{The complete positive-cone restriction and its probability law}

Let $J_{++}$ denote the original positive definite cone in $J$. Define
the following explicit open subset of the metric target:
\begin{equation}\label{rm:positive-image}
 \mathcal B_+=\{y\in\mathcal B:
        a>0,\ \Delta>0,\ m(y)+\beta(y)>0\}.
\end{equation}

\begin{theorem}\label{rm:positive-theorem}
The exact image of $\mathcal U\cap J_{++}$ under $\Psi$ is
$\mathcal B_+$. For a fixed fibre, all its points are positive definite
if and only if $a>0$, $\Delta>0$, and $m>\beta$.
For $a>0$, $\Delta>0$, the Haar probability of positive definiteness is
\begin{equation}\label{rm:positive-probability}
 p(y)=
 \begin{cases}
  0,&m\leq-\beta,\\
  1-\pi^{-1}\arccos(m/\beta),&-\beta<m<\beta,\\
  1,&m\geq\beta.
 \end{cases}
\end{equation}
If $a\leq0$ or $\Delta\leq0$, this probability is zero regardless of
$m$. At $m=\beta$ the fibre has singular positive semidefinite boundary points of Haar
measure zero, although $p=1$. On $\mathcal B_+$ the probability kernel
\begin{equation}\label{rm:positive-kernel}
 \kappa_y^+(E)=p(y)^{-1}\kappa_y(E\cap J_{++})
\end{equation}
is well defined and its cubic law is precisely
\begin{equation}\label{rm:positive-law}
 N_*\kappa_y^+(dt)
 =\frac{\mathbf1_{(0,\infty)\cap(m-\beta,m+\beta)}(t)}
       {\pi p(y)\sqrt{\beta^2-(t-m)^2}}\,dt.
\end{equation}
For $-\beta<m<\beta$ put $\gamma=\arccos(m/\beta)$ and
$s=\sin\gamma$. Its first two moments and variance are
\begin{align}
 \mathbb E_{\kappa_y^+}N
   &=m+\frac{\beta s}{\pi p},\label{rm:positive-mean}\\
 \mathbb E_{\kappa_y^+}N^2
   &=m^2+\frac{\beta^2}{2}
                    +\frac{3m\beta s}{2\pi p},\label{rm:positive-second}\\
 \operatorname{Var}_{\kappa_y^+}N
   &=\frac{\beta^2}{2}-\frac{m\beta s}{2\pi p}
                     -\frac{\beta^2s^2}{\pi^2p^2}.
                     \label{rm:positive-variance}
\end{align}
For $m\geq\beta$ these quantities equal the unrestricted values
$m$, $m^2+\beta^2/2$, and $\beta^2/2$, respectively.
\end{theorem}
\begin{proof}
The upper block is
$K=\left(\begin{smallmatrix}a&v\\\bar v&b\end{smallmatrix}\right)$
with $v=c+\bar\rho\mathbf j$, so
$\Delta=ab-|v|^2$. For $a>0$, completion of the square gives
\[
 (x_1,x_2)^*K(x_1,x_2)
 =a|x_1+a^{-1}v x_2|^2+(b-|v|^2/a)|x_2|^2.
\]
It follows that $K>0$ if and only if $a>0$ and $\Delta>0$.
Necessity of $a>0$ follows by testing $(1,0)$, so this covers every
case. For such $K$ its inverse is exactly
\[
 K^{-1}=\Delta^{-1}\begin{pmatrix}b&-v\\-\bar v&a\end{pmatrix},
\]
as multiplication verifies, since $a,b$ are real and
$v\bar v=\bar vv=|v|^2$. For a quaternionic column $x$ and scalar $t$,
\[
 (x,t)^*M(x,t)
  =(x+K^{-1}qt)^*K(x+K^{-1}qt)
      +(\lambda-q^*K^{-1}q)|t|^2.
\]
This proves $M>0$ exactly when $K>0$ and
$\lambda-q^*K^{-1}q>0$. Expanding the displayed inverse shows
\[
 \Delta(\lambda-q^*K^{-1}q)
 =\lambda\Delta-b|q_1|^2-a|q_2|^2
                       +2\operatorname{Re}(\bar q_1vq_2)=N(M).
\]
Consequently, on every fibre with $a>0,\Delta>0$, positivity is
equivalent to the scalar condition $N>0$. Congruence by the original
invertible $X$ preserves positivity because
$v^*XMX^*v=(X^*v)^*M(X^*v)$ and $X^*$ is bijective.

The phase formula attains every value in $[m-\beta,m+\beta]$.
There is at least one strictly positive point precisely when
$m+\beta>0$, and every point is strictly positive precisely when
$m-\beta>0$. This proves the exact image and full-fibre statements.
All inequalities in \eqref{rm:positive-image} are strict and involve
continuous functions, proving openness directly.

With the uniform variable $t$ used in the proof of
\eqref{rm:arcsine-law}, positivity is $\cos t<m/\beta$.
For $-\beta<m<\beta$ this means $\gamma<t<2\pi-\gamma$, which
has probability $1-\gamma/\pi$. The endpoint and outer cases give
\eqref{rm:positive-probability}. Restricting and dividing the
arcsine law by this positive mass gives
\eqref{rm:positive-kernel}--\eqref{rm:positive-law}, including the
boundary case $m=\beta$ where only zero-probability points are
removed. The measure is Borel by the kernel property proved above.

Finally direct integrals over this interval give
\[
 \frac1{2\pi}\int_\gamma^{2\pi-\gamma}\cos t\,dt
       =-\frac{s}{\pi},\qquad
 \frac1{2\pi}\int_\gamma^{2\pi-\gamma}\cos^2t\,dt
       =\frac p2-\frac{\sin(2\gamma)}{4\pi}.
\]
Insert $N=m-\beta\cos t$, divide by $p$, and use
$\cos\gamma=m/\beta$ and $\sin(2\gamma)=2s m/\beta$.
This yields \eqref{rm:positive-mean} and
\eqref{rm:positive-second}. Subtracting the square of the mean
gives \eqref{rm:positive-variance}. For $m\geq\beta$ the excluded
phase set has zero probability, so the unrestricted moments apply.
\end{proof}

This theorem identifies exactly why a positivity restriction cannot
always retain uniform probability on the entire original fibre.
The group still determines the unrestricted measure; the chosen state
space $J_{++}$ requires the explicit restriction
\eqref{rm:positive-kernel}. Its four discrete components retain equal
mass because the right $SU(2)$ congruence preserves positivity and the
phase law is identical on each component.

There are regular fibres crossing the boundary of the positive cone,
so this distinction occurs in the actual construction. One exact
choice, avoiding numerical approximations to logarithms, is
\begin{equation}\label{rm:crossing-example}
 G=4I_2,\quad S=\operatorname{diag}(2,1),\quad
 \rho=\frac12+i,\quad\lambda=\frac{80}{59}.
\end{equation}
Here $\Delta=16-5/4=59/4$, $T=\operatorname{diag}(4,1)$,
$m=(80/59)(59/4)-20=0$, and $\beta=2\sqrt5$.
The weights are $(2,e^{1/2},e,e^{80/59})$.
They have $d_1<d_2<d_3$ because $1/2<1<80/59$;
$e^{1/2}<2$ follows from $e<4$, and $2<e$ follows from
$e=\sum_{n\geq0}1/n!>2$. To check the former bound explicitly,
$n!\geq2^{n-1}$ for $n\geq2$, with a strict inequality for some $n$,
so $e<2+\sum_{n\geq2}2^{-(n-1)}=3<4$.
Thus all four weights are distinct. We have $a>0$ and $\Delta>0$,
but half of each fibre circle is positive and half negative in the
Schur complement. On the positive half,
\[
 p=\frac12,\qquad
 \mathbb E N=\frac{4\sqrt5}{\pi},\qquad
 \operatorname{Var}N=10-\frac{80}{\pi^2}.
\]
The matrix at phase zero has $N=-2$ and the one at phase $\theta=\pi/2$
has $N=2$, by \eqref{rm:cubic-fourier}. These are two explicitly
retained matrices in the same regular metric fibre with opposite
positivity status, while the upper block remains positive.

\subsection{An explicit weighted family and the information it retains}

The compact group supplies an invariant conditional probability. A
different conditional weight is an additional model choice whose
effect can also be calculated without selecting a frame. For each
fixed $\epsilon\in[-1,1]$, define the globally specified weight
\begin{equation}\label{rm:weighted-family}
 w_\epsilon(H)=1+\epsilon
       \frac{N(H)-m(\Psi(H))}{\beta(\Psi(H))},\qquad
 d\kappa_y^{(\epsilon)}(H)=w_\epsilon(H)\,d\kappa_y(H).
\end{equation}
Because $|N-m|\leq\beta$ and $\beta>0$ on the regular domain,
this is a nonnegative Borel weight. Its integral is one by
\eqref{rm:first-two-moments}, so it gives an actual probability
kernel with no hidden normalizing constant. It does not depend on
which representative of the fibre was used to construct $\kappa_y$.
The centred variable $Z_y=N-m$ has symmetric law, with
$\mathbb E Z_y=\mathbb E Z_y^3=0$ and
$\mathbb E Z_y^2=\beta^2/2$. Therefore
\begin{equation}\label{rm:weighted-moments}
 \mathbb E_{\kappa_y^{(\epsilon)}}N=m+\frac{\epsilon\beta}{2},
 \qquad
 \operatorname{Var}_{\kappa_y^{(\epsilon)}}N
       =\frac{\beta^2}{2}-\frac{\epsilon^2\beta^2}{4}.
\end{equation}
Indeed weighting multiplies each expectation by
$1+\epsilon Z_y/\beta$; it changes the mean of $Z_y$ to
$\epsilon\beta/2$ and leaves its second moment equal to
$\beta^2/2$, which proves both formulas.

For $\epsilon\ne0$ this weight need not be invariant under the scalar
phase action: the displayed cubic has a nonconstant phase mode.
It remains a well-defined target-indexed conditional distribution.
This explicitly distinguishes the compact-group invariant choice
from a specified weighted choice by their measurable predictions,
including the mean shift $\epsilon\beta/2$.

The exact map supplied by this construction is thus
\[
 J\supset\mathcal U\xrightarrow{\Psi}\mathcal B,
 \qquad y\longmapsto\kappa_y,
 \qquad f\longmapsto\mathcal A f,
 \qquad y\longmapsto N_*\kappa_y.
\]
For the original cubic the last map is explicitly
\eqref{rm:arcsine-law}. It retains the complete conditional cubic
distribution, including the positive-cone support and fluctuation
correction \eqref{rm:cubic-penalty}. Strictly positive fibre variance
proves that no deterministic scalar $n(y)$ recovers $N(H)$ on an
entire regular fibre. The probability kernel is the exact relation
that remains: it transports observables by conditional integration,
and transports the cubic to a fully computed distribution.
